% Non-Monotonic Optimum in Prefix-Token Perturbation of Language Model Reasoning
% NeurIPS 2026 Workshop Submission — DRAFT
%
% Sections marked [LOCKED] are stable and can be finalized.
% Sections marked [PENDING] await experimental data.
% Sections marked [PLACEHOLDER] need substantial revision.

\documentclass{article}
\usepackage[utf8]{inputenc}
\usepackage{amsmath,amssymb}
\usepackage{booktabs}
\usepackage{graphicx}
\usepackage{hyperref}
\usepackage{xcolor}
\usepackage{subcaption}

% NeurIPS workshop format (if needed, swap for neurips_2026.sty)
\usepackage[margin=1in]{geometry}

\newcommand{\pending}[1]{\textcolor{red}{[PENDING: #1]}}
\newcommand{\placeholder}[1]{\textcolor{orange}{[PLACEHOLDER: #1]}}

\title{Random Prefix Perturbations Diversify Reasoning Trajectories\\with a Sharp Non-Monotonic Optimum in Small Language Models}

\author{Devansh}
\date{March 2026}

\begin{document}
\maketitle

% ============================================================================
\begin{abstract}
Prepending random embedding-scale tokens to the input of a 4-bit quantized
language model (Qwen3-4B) produces a sharp, non-monotonic improvement in
chain-of-thought arithmetic reasoning. Two random tokens raise mean
accuracy from 32\% to 51.6\% (+19.6 percentage points) across ten
independent directions, while one token yields 42.7\% and three or eight
tokens revert to $\sim$44\%. A first-token logit probe confirms that
perturbation does not activate reasoning mode (which is already
saturated at $>$99.99\% under thinking-enabled decoding); instead,
it modulates the downstream reasoning trajectory. Different
perturbation directions solve \emph{different task subsets}: ten two-token
directions collectively cover 100\% of tasks (25/25), compared to 80\%
for three-token and 92\% for eight-token perturbations. The three-token
oracle set is a strict subset of the two-token oracle, indicating that
over-perturbation contracts the reachable task space. Cross-model validation confirms the effect on four models including
Qwen3-8B at 8-bit (+12.8pp at $n{=}10$, oracle 80\%,
McNemar 16/0), DeepSeek-R1-Distill-1.5B (oracle 100\%,
McNemar 6/0), and phi-2 (+6.7pp, out-of-family). A within-model
quantization control reveals that weight precision modulates
perturbation sensitivity: Qwen3-8B at 4-bit yields a null result
while 8-bit yields +16pp on the same tasks. DeepSeek exhibits
task-selective recoverability: every baseline failure is rescued by at
least one perturbation direction, though the average random direction
is slightly harmful ($-1.6$pp at $n{=}10$), establishing a
distinction between oracle accessibility and mean-effect improvement.
A dose-response scout on DeepSeek confirms the non-monotonic window
generalizes: 1-token perturbation hurts ($-12$pp), 2-token is the
local optimum, and 3-token enters a direction-sensitive bifurcated
regime. A grading-confound analysis reveals that the accuracy gains are
primarily a \emph{convergence} effect: the model can already compute
correct answers 80\% of the time (answer-anywhere), but perturbation
improves the rate at which the correct answer appears as the final
output (32\% $\to$ 51.6\%).
Cross-domain validation on five complex planning tasks reveals two
additional phenomena: (1)~perturbation rescues catastrophic
attention-sink-induced generation failures (14-word baseline collapse
$\to$ 650+ word complete plans across all five perturbation seeds),
and (2)~\emph{evolved} soft prompts (via a trained latent scorer +
evolutionary search) surface qualitatively different reasoning---security
frameworks, architectural patterns, and diagnostic strategies absent
from both baseline and random perturbation outputs.
Further cross-domain validation on twelve complex legal reasoning tasks
(blind-reviewed by an independent LLM judge) confirms that the best-of-5
perturbation output beats greedy baseline in 11/12 tasks (average
lift $+1.6$ on a 10-point expert-review scale). On tasks where greedy
decoding produces truncated or legally flawed analysis, perturbation
unlocks complete multi-jurisdictional reasoning the model possesses but
cannot access under standard generation. These results confirm that the
system is explicitly \emph{judge-heavy}: the perturbation mechanism
reliably accesses latent knowledge, and scorer/judge quality determines
how much of that ceiling is captured.
\end{abstract}

% ============================================================================
\section{Introduction} % [LOCKED — stable claims only]

Large language models possess reasoning capabilities that are not
uniformly accessible through standard prompting. Recent work has shown
that semantically meaningless prefix tokens---repeated punctuation
\cite{shi2025}, pause tokens \cite{goyal2024}, and trailing whitespace
\cite{salinas2024}---can measurably alter model behavior, sometimes
improving reasoning performance.

We study a related but distinct regime: \emph{continuous} perturbation
in embedding space. Rather than prepending discrete tokens from the
vocabulary, we inject random vectors at the model's native embedding
scale as soft prompt prefixes. This eliminates semantic content entirely,
isolating the effect of perturbation energy and position from token
identity.

Our key findings on Qwen3-4B (4-bit quantized) with 25 nested arithmetic
tasks are:

\begin{enumerate}
    \item \textbf{Sharp non-monotonic optimum.} Two random prefix tokens
    raise mean last-integer accuracy from 32\% to 51.6\% across ten
    directions (primarily a convergence effect; answer-anywhere accuracy
    is 80\% at baseline); one token yields 42.7\%, and three or eight
    tokens revert to $\sim$44\%.

    \item \textbf{Complete oracle coverage.} Ten two-token directions
    collectively solve 100\% of tasks (25/25), compared to 80\% (20/25)
    for three-token and 92\% (23/25) for eight-token perturbations.
    The three-token oracle set is a strict subset of the two-token
    oracle: five tasks solvable at two tokens become frozen at three.
    At matched budget ($k = 3$), two tokens achieve 77\% oracle coverage
    versus 60\% for three tokens and 68\% for eight tokens.

    \item \textbf{Heterogeneous task selectivity.}
    Different perturbation directions solve different task subsets,
    but per-direction solve counts vary normally (std = 1.87,
    $p = 0.66$ vs.\ heterogeneous iid null). The diversity is in
    \emph{which} tasks are solved, not in \emph{how many}.

    \item \textbf{Trajectory modulation, not mode gating.}
    A first-token logit probe shows that \texttt{<think>} probability
    is saturated ($>$99.99\%) under all conditions including baseline.
    Perturbation does not activate reasoning mode; it modulates the
    downstream reasoning trajectory within an already-active think
    mode.

    \item \textbf{Cross-domain: attention sink avoidance.}
    On complex planning tasks, greedy baseline can fail catastrophically
    (14 words before collapsing into degenerate output). Two-token
    perturbation rescues these failures into complete 650+ word plans
    across all five seeds, demonstrating that soft prompt perturbation
    breaks attention-sink-induced degenerate generation paths.

    \item \textbf{Cross-domain: latent knowledge access via evolution.}
    Evolved soft prompts (trained scorer + evolutionary search) surface
    qualitatively different reasoning---security frameworks (MITRE
    ATT\&CK, HSM-backed credential rotation), architectural patterns
    (honeypot deployment, DMZ isolation), and diagnostic strategies
    (app-level idempotency analysis) absent from both baseline and
    random perturbation outputs. This suggests embedding-space
    perturbation can access latent knowledge the model possesses but
    cannot reach under standard decoding.

    \item \textbf{Cross-domain: legal reasoning validated by blind review.}
    On twelve complex legal reasoning tasks (FTC regulation, employment
    discrimination, IP risk, contract review, negotiation strategy),
    the best-of-5 perturbation output beats greedy baseline in 11/12
    tasks (average $+1.6$ on 10-point blind expert-review
    scale). Perturbation accesses multi-jurisdictional legal analysis
    the model cannot reach via greedy decoding. The system is explicitly
    judge-heavy: scorer quality determines outcome quality.
\end{enumerate}

All experiments use greedy decoding (temperature = 0), making the
observed variability across perturbation directions a form of
\emph{deterministic chaos} rather than sampling noise.

% ============================================================================
\section{Related Work} % [LOCKED — core refs stable]

\paragraph{Discrete prefix perturbations.}
Shi et al.~\cite{shi2025} demonstrate that repeated punctuation tokens
(e.g., ``/'' and ``?'') improve reasoning accuracy by 1--5\% across
multiple models, attributing the effect to MLP activation redistribution.
Our work studies the continuous embedding-space regime, finding
substantially larger effects (+19.6pp mean at the optimum), direction
independence, and task-selective redistribution---phenomena not observed
in discrete token space. In a head-to-head comparison on our task set
(Section~\ref{sec:shi}), discrete ``/'' and ``?'' tokens at 2 tokens
yield +4pp and +16pp respectively (mean +10pp), while continuous random
noise yields +19.6pp---a 9.6pp advantage for continuous perturbation.

\paragraph{Pause tokens.}
Goyal et al.~\cite{goyal2024} show that training with learned pause
tokens improves downstream performance. Our effect requires no training:
random embeddings at inference time suffice.

\paragraph{Extra-token expressivity.}
London \& Nagarajan~\cite{london2025} prove that additional tokens
increase transformer expressivity. Our results provide an empirical
instantiation: two tokens at the correct energy scale unlock reasoning
that the model cannot access with standard prompting.

\paragraph{Deterministic chaos in transformers.}
Li et al.~\cite{li2025} introduce quasi-Lyapunov exponents for
language models, establishing that small perturbations to hidden states
grow exponentially through transformer layers. Our invariance-length
analysis (Section~\ref{sec:chaos}) provides a complementary behavioral
measure of this phenomenon under greedy decoding.

\paragraph{Attention sinks.}
Xiao et al.~\cite{xiao2024} document the attention sink phenomenon:
early tokens accumulate disproportionate attention from all subsequent
tokens regardless of semantic content. Our cross-domain planning
results (Section~\ref{sec:planning}) provide a practical consequence:
when attention sink positions contain only standard prompt tokens,
greedy decoding can lock into degenerate generation paths. Soft prompt
perturbation in these positions disrupts the degenerate pattern,
rescuing catastrophic failures.

\paragraph{Stochastic resonance.}
The non-monotonic dose-response is reminiscent of stochastic resonance
\cite{gammaitoni1998}, where noise improves signal detection in nonlinear
systems. However, our statistical tests do not confirm the formal
inverted-U relationship at the task level (Section~\ref{sec:limitations}),
so we use the more conservative framing of ``non-monotonic optimum.''

\paragraph{Continuous embedding perturbation.}
Concurrent and independent work confirms our central observation.
Kim et al.~\cite{kim2026rsp} show that Random Soft Prompts---untrained
embedding vectors sampled from the pretrained distribution---achieve
accuracy comparable to trained soft prompts on math reasoning benchmarks.
Zhu et al.~\cite{zhu2025softreasoning} combine first-token embedding
perturbation with Bayesian optimization and a trained verifier, achieving
strong results on GSM8K and related tasks (ICML 2025 Spotlight). Our work
differs in three ways: (1)~we use deterministic greedy decoding to isolate
perturbation effects from sampling noise, (2)~we study the
\emph{selector realization problem}---how much oracle headroom can be
recovered without external labels or verifiers, and (3)~we demonstrate
multi-domain generalization beyond arithmetic.

\paragraph{Hidden-state reward models.}
Guo et al.~\cite{guo2025swift} show that simple linear probes on LLM
hidden states can score best-of-$N$ candidates more accurately than
dedicated reward models, using less than 0.005\% of the generating
model's parameters. Our selector heuristics operate on candidate-level
output features; combining them with hidden-state probes is a promising
direction for selector improvement.

% ============================================================================
\section{Methods} % [LOCKED]

\subsection{Model and Inference}

We use Qwen3-4B~\cite{qwen2025} with 4-bit GPTQ quantization,
running on a single NVIDIA RTX 5090 Laptop GPU (24GB VRAM). All
generation uses greedy decoding (temperature = 0, \texttt{do\_sample =
False}) with \texttt{max\_new\_tokens = 1024} and thinking mode enabled.

\subsection{Task Suite}

We use 25 nested arithmetic expressions at ``sweet-spot'' difficulty
(3-step expressions with intermediate results in the hundreds to
thousands range, 32\% baseline accuracy). Tasks are deterministically
generated and fixed across all conditions.

\subsection{Perturbation Protocol}

We inject $k$ random embedding vectors as soft prompt tokens before
the user input. Each vector is sampled from $\mathcal{N}(0, I)$ and
scaled to match the model's embedding table RMS norm
($\sigma_{\text{emb}} \approx 0.022$), calibrated via
\texttt{auto\_calibrate()} over 1000 vocabulary samples.

\subsection{Conditions}

\begin{table}[h]
\centering
\small
\begin{tabular}{llp{5cm}}
\toprule
\textbf{Condition} & \textbf{Tokens} & \textbf{Description} \\
\midrule
Baseline & 0 & No prefix tokens \\
Zero embedding & 8 & All-zero vectors \\
Mean embedding & 8 & Mean of embedding table \\
Force-think & 0 & Forced \texttt{<think>} token, no noise \\
Random noise & 1, 2, 3, 8 & Random $\mathcal{N}(0,I)$, RMS-matched \\
Discrete tokens & 2 & Repeated ``/'' and ``?'' tokens (Shi comparison) \\
\bottomrule
\end{tabular}
\caption{Experimental conditions. All use greedy decoding.}
\label{tab:conditions}
\end{table}

\subsection{Scoring and Data Integrity}

Correctness is determined by \texttt{verify\_answer()}: the last integer
in the generated response must exactly match the expected answer. All
experiments with the current codebase store the extracted answer alongside
the full response (up to 2000 characters) for independent verification.

% ============================================================================
\section{Theoretical Framework: Basin Model and Selector Gap}
\label{sec:theory}

We formalize the relationship between perturbation diversity and
realized accuracy using a \emph{basin model} of deterministic
trajectory spaces.

\subsection{Perturbation as Deterministic Function}

Let $M$ be a frozen language model with greedy decoding and let $p$
denote a prompt. A latent perturbation $z \in \mathbb{R}^d$ defines a
soft prompt prefix that, under greedy decoding, produces a deterministic
output:
\[
y = M(p, z).
\]
Since decoding is deterministic, the map $z \mapsto M(p, z)$ partitions
the latent space into \emph{basins}: regions (not necessarily connected)
that produce identical or equivalent outputs.

\subsection{Oracle Coverage and Hit Probability}

For a task with prompt $p$ and verifier $V$, define the \emph{success
set} $S_p = \{z : V(M(p, z)) = 1\}$. Drawing $k$ perturbations
$z_1, \ldots, z_k$ independently, the \emph{oracle accuracy} is:
\[
\text{Acc}_{\text{oracle}}(k) = \Pr\!\left[\exists\, i : z_i \in S_p\right]
\approx 1 - (1 - \hat{p})^{k_{\text{eff}}},
\]
where $\hat{p}$ is the per-candidate hit probability and $k_{\text{eff}}
= k / (1 + (k-1)\rho)$ is a design-effect approximation that adjusts
for pairwise correlation $\rho$ among candidates. When candidates are
independent ($\rho = 0$), $k_{\text{eff}} = k$; when perfectly
correlated ($\rho = 1$), $k_{\text{eff}} = 1$ regardless of $k$.
This is an approximation; the exact oracle probability under correlated
Bernoulli hits depends on the full correlation structure.

Fitting $\log(1 - \text{Acc}_{\text{oracle}})$ versus $k$ estimates
$\hat{p}$ and reveals whether additional perturbations provide
diminishing returns (high $\rho$) or near-independent coverage
(low $\rho$).

\subsection{The Selector Gap}

The oracle provides an upper bound on what perturbation diversity
\emph{could} achieve, but it requires ground-truth labels. In
deployment, a \emph{selector} $\sigma$ must choose among candidates
using only local features (token counts, answer consistency, response
structure). We define:

\begin{itemize}
    \item \textbf{Oracle recovery:} $R_\sigma = \frac{\text{Acc}_\sigma - \text{Acc}_{\text{majority}}}{\text{Acc}_{\text{oracle}} - \text{Acc}_{\text{majority}}}$
    \item \textbf{Regret:} $\text{Acc}_{\text{oracle}} - \text{Acc}_\sigma$
    \item \textbf{Selector gap:} The difference between oracle coverage and the best realized label-free accuracy.
\end{itemize}

The selector gap is the central scientific object: it quantifies how
many of the successful trajectories reachable under this perturbation
scheme and verifier are \emph{accessible} (oracle) versus
\emph{realizable} (selector) without external supervision. A narrow
gap means label-free deployment is viable; a wide gap motivates better
selector design.

\subsection{Effective Diversity}

Not all $k$ candidates contribute independent information. We measure
\emph{effective diversity} as the normalized entropy of the answer
distribution:
\[
D_{\text{eff}} = \frac{H(\text{answers})}{\log_2 k},
\]
where $H$ is the Shannon entropy over extracted answer clusters.
$D_{\text{eff}} = 0$ when all candidates agree (no diversity);
$D_{\text{eff}} = 1$ when all candidates give distinct answers
(maximum diversity). Note that high diversity is neither necessary
nor sufficient for oracle coverage: all candidates may agree on
the correct answer ($D_{\text{eff}} = 0$, oracle hit), or produce
many distinct wrong answers ($D_{\text{eff}} \approx 1$, oracle miss).
However, across tasks, effective diversity indicates the degree to
which perturbation explores distinct trajectories.

\subsection{Task Taxonomy}

The basin model induces a natural task taxonomy based on the success
set $S_p$:
\begin{itemize}
    \item \textbf{Unanimous:} $\hat{p} \approx 1$; greedy already solves the task, and all perturbations agree. No selector needed.
    \item \textbf{Sensitive:} $0 < \hat{p} < 1$; some perturbations solve the task, others do not. This is the regime where selectors matter.
    \item \textbf{Frozen:} $\hat{p} \approx 0$; no perturbation reaches the success set. The task requires capabilities beyond what perturbation can access.
\end{itemize}

The majority failure rate---tasks where the oracle succeeds but
majority vote fails---directly measures the selector gap for the
majority-vote selector and motivates the search for better
label-free selectors.

% ============================================================================
\section{Results} % [LOCKED sections noted; PENDING sections noted]

\subsection{Non-Monotonic Dose-Response} % [LOCKED]
\label{sec:dose}

\begin{table}[h]
\centering
\begin{tabular}{cccccc}
\toprule
\textbf{Tokens} & \textbf{Accuracy} & \textbf{Std} & \textbf{$n$} & \textbf{$\Delta$ vs baseline} \\
\midrule
0 & 32.0\% & --- & 1 & --- \\
1 & 42.7\% & 1.9\% & 3 & +10.7pp \\
\textbf{2} & \textbf{51.6\%} & \textbf{7.9\%} & \textbf{10} & \textbf{+19.6pp} \\
3 & 44.0\% & 5.3\% & 10 & +12.0pp \\
8 & 44.4\% & 7.0\% & 10 & +12.4pp \\
\bottomrule
\end{tabular}
\caption{Dose-response: accuracy vs.\ number of random prefix tokens
($n$ = number of independent directions tested).
Two tokens is the sharp optimum. Three and eight tokens revert to
$\sim$44\%.}
\label{tab:dose}
\end{table}

\begin{figure}[h]
\centering
\includegraphics[width=0.7\linewidth]{experiments/fig_dose_response.pdf}
\caption{Non-monotonic dose-response curve. Error bars show standard
error of the mean across $n$ perturbation directions.}
\label{fig:dose}
\end{figure}

The relationship between token count and accuracy is sharply non-monotonic.
Two random tokens produce the highest mean accuracy (51.6\% $\pm$ 7.9\%
across ten independent directions). Adding a third token drops accuracy to
$\sim$44\%, matching the eight-token level. The three-token oracle set
(20/25) is a strict subset of the two-token oracle (25/25), indicating
that over-perturbation contracts the reachable task space rather than
merely redistributing it.

\subsection{Explicit Think-Prefix Control} % [LOCKED]
\label{sec:decomposition}

\begin{table}[h]
\centering
\begin{tabular}{lcc}
\toprule
\textbf{Condition} & \textbf{Accuracy} & \textbf{$\Delta$ vs baseline} \\
\midrule
Baseline & 32.0\% & --- \\
Explicit \texttt{<think>} prefix & 40.0\% & +8pp \\
2-tok random noise (mean, $n = 10$) & 51.6\% & +19.6pp \\
\bottomrule
\end{tabular}
\caption{Explicit think-prefix control. Teacher-forcing the first token
to \texttt{<think>} adds +8pp. Random noise adds a further +11.6pp
beyond the explicit prefix (mean across 10 directions). Note: a
first-token logit probe (Section~\ref{sec:probe}) shows that
\texttt{<think>} probability is already $>$99.99\% at baseline under
thinking-enabled decoding, so the +8pp is not mode activation but a
trajectory-priming effect of the explicit prefix.}
\label{tab:decomposition}
\end{table}

A first-token logit probe (Section~\ref{sec:probe}) confirms that
\texttt{<think>} is the dominant first token under all conditions
(probability $>$99.99\%), including unperturbed baseline. The observed
accuracy differences therefore reflect trajectory modulation within
an already-active reasoning mode, not mode-switching. Teacher-forcing
\texttt{<think>} as an explicit prefix adds +8pp; random noise provides
an additional +11.6pp beyond this, indicating that the perturbation
effect operates on the reasoning trajectory rather than on mode entry.

\subsection{Oracle Efficiency} % [LOCKED]
\label{sec:oracle}

Because different perturbation directions solve different task subsets,
running multiple perturbations and taking the oracle (any-correct)
provides coverage well above individual accuracy. Figure~\ref{fig:oracle}
shows coverage as a function of budget $k$ for each condition.

\begin{figure}[h]
\centering
\includegraphics[width=0.8\linewidth]{experiments/fig_oracle_coverage.pdf}
\caption{Oracle coverage vs.\ perturbation budget. Two-token (green)
reaches 100\% coverage at $k = 10$ (zero frozen tasks), dominates at
every matched budget, and its oracle set is a strict superset of the
three-token set (red, 80\% at $k = 10$). Eight-token (blue) reaches
92\% at $k = 10$.}
\label{fig:oracle}
\end{figure}

The non-monotonic optimum extends beyond per-direction accuracy to
coverage efficiency: five tasks solvable at two tokens become frozen
at three tokens, contracting the reachable task set rather than merely
redistributing it. Two-token perturbation alone covers all 25 tasks;
combined oracle across all conditions is therefore 25/25 = 100\%.

Oracle coverage grows more slowly than the independence null (positive
task correlation, $z = -1.35$ for 2-tok vs.\ $z = -8.25$ for 8-tok),
indicating that latents fail on correlated hard tasks. The two-token
condition has the least redundancy, closest to independence.

\subsection{Task Selectivity Without Equalization}
\label{sec:equalization}

An initial observation at $n = 3$ two-token directions suggested
solve-count equalization (all directions solving exactly 15/25 tasks,
$p = 0.031$ vs.\ heterogeneous iid null). This did not replicate:
at $n = 10$ directions, solve counts range from 9 to 15
(std = 1.87, $p = 0.66$), fully consistent with independent sampling
at heterogeneous per-task solve rates. The initial pattern was
small-sample noise. The robust finding is \emph{task selectivity}:
directions solve different subsets, enabling oracle gains, but the
total number solved per direction varies normally.

We test equalization by Monte Carlo simulation: given the observed
per-task solve rates, we simulate 100{,}000 draws under a
heterogeneous iid Bernoulli null (each task independently solved at
its observed rate) and compare the observed standard deviation of
per-direction solve counts against this null distribution.

We decompose tasks at each perturbation level into three categories:
\emph{unanimous} (solved by all sampled directions), \emph{frozen}
(solved by none), and \emph{sensitive} (solved by some but not all).
These categories are probabilistic, not absolute---additional directions
may solve ``frozen'' tasks or fail ``unanimous'' ones.

\begin{table}[h]
\centering
\textbf{Panel A: Matched budget ($k = 3$ directions)}

\medskip
\begin{tabular}{lcccc}
\toprule
\textbf{Condition} & \textbf{Unan.} & \textbf{Frozen} & \textbf{Sens.} & \textbf{Mean acc.} \\
\midrule
1-tok & 6.0 & 8.0 & 11.0 & 10.7/25 \\
\textbf{2-tok}$^\dagger$ & \textbf{7.6} & \textbf{5.8} & \textbf{11.6} & \textbf{12.9/25} \\
3-tok$^\dagger$ & 7.3 & 9.9 & 7.8 & 11.0/25 \\
8-tok$^\dagger$ & 6.0 & 7.9 & 11.2 & 11.1/25 \\
\bottomrule
\end{tabular}

\medskip
\textbf{Panel B: Full observed directions (with iid $p$-values)}

\medskip
\begin{tabular}{lccccccc}
\toprule
\textbf{Condition} & $k$ & \textbf{Unan.} & \textbf{Frozen} & \textbf{Sens.} & \textbf{Mean/25} & \textbf{Std} & $p_{\text{iid}}$ \\
\midrule
1-tok & 3 & 6 & 8 & 11 & 10.7 & 0.47 & 0.23 \\
\textbf{2-tok} & \textbf{10} & \textbf{6} & \textbf{0} & \textbf{19} & \textbf{12.9} & \textbf{1.87} & \textbf{0.659} \\
3-tok & 10 & 5 & 5 & 15 & 11.0 & 1.33 & 0.335 \\
8-tok & 10 & 3 & 2 & 20 & 11.1 & 1.76 & 0.586 \\
\bottomrule
\end{tabular}
\caption{Task decomposition by perturbation magnitude.
\textbf{Panel A}: At matched budget ($k = 3$), two tokens minimize the
frozen set and maximize per-direction accuracy.
$^\dagger$Values are averaged over all $\binom{n}{3}$ subsets
(120 combinations each for $n = 10$ conditions).
\textbf{Panel B}: Full observed directions with Monte Carlo
$p$-values against heterogeneous iid null. No condition shows
significant variance suppression; the diversity is in task identity,
not solve count.}
\label{tab:equalization}
\end{table}

\subsection{Deterministic Chaos} % [LOCKED]
\label{sec:chaos}

Under greedy decoding, different random prefix vectors produce
byte-identical text for an ``invariance length'' that decreases with
perturbation energy:

\begin{table}[h]
\centering
\begin{tabular}{ccc}
\toprule
\textbf{Tokens} & \textbf{Mean Invariance} & \textbf{Oracle} \\
\midrule
1 & 349 chars & 68\% \\
2 & 323 chars & 100\% \\
8 & 93 chars & 92\% \\
\bottomrule
\end{tabular}
\caption{Invariance length (byte-identical prefix before divergence)
under greedy decoding. This is deterministic chaos, not sampling noise.}
\label{tab:invariance}
\end{table}

The two-token sweet spot balances long enough invariance to preserve
reasoning quality with short enough invariance for task-subset diversity.

\subsection{Controls} % [LOCKED]

\begin{itemize}
    \item \textbf{Zero embedding} (+4pp): Sequence extension alone
    provides minimal benefit.
    \item \textbf{Mean embedding} (+4pp): Identical to zero---no
    advantage from the ``average token.''
    \item \textbf{No-think mode}: When thinking is disabled
    (\texttt{enable\_thinking=False}), perturbation has no effect.
    Chain-of-thought is the mediating mechanism.
    \item \textbf{Timing confound resolved}: The highest-accuracy
    2-tok directions achieve 60\% (15/25) at baseline-identical mean
    timing (73.6s), ruling out extra computation as the mechanism.
\end{itemize}

\subsection{First-Token Logit Probe} % [LOCKED]
\label{sec:probe}

To test whether perturbation activates reasoning mode, we extract
logits at the first decode position (immediately after the assistant
turn) under all conditions. The \texttt{<think>} token is rank 1
with probability $>$99.99\% for all 25 tasks under every condition,
including unperturbed baseline (Table~\ref{tab:probe}).

\begin{table}[h]
\centering
\small
\begin{tabular}{lcc}
\toprule
\textbf{Condition} & \textbf{Mean $P(\texttt{<think>})$} & \textbf{Rank = 1} \\
\midrule
Baseline & 99.994\% & 25/25 \\
Zero (2-tok) & 99.994\% & 25/25 \\
Mean (2-tok) & 99.994\% & 25/25 \\
Noise 1-tok & 99.996\% & 25/25 \\
Noise 2-tok & 99.997\% & 25/25 \\
Noise 3-tok & 99.997\% & 25/25 \\
Noise 8-tok & 99.995\% & 25/25 \\
Discrete \texttt{/} (2-tok) & 99.998\% & 25/25 \\
Discrete \texttt{?} (2-tok) & 99.997\% & 25/25 \\
\bottomrule
\end{tabular}
\caption{First-token logit probe. \texttt{<think>} probability is
saturated under all conditions including baseline. Perturbation does
not activate reasoning mode; it modulates the subsequent reasoning
trajectory.}
\label{tab:probe}
\end{table}

This rules out first-token mode gating as the mechanism. The
perturbation effect operates within an already-active reasoning mode,
modulating the trajectory of the chain-of-thought rather than
determining whether the model reasons at all.

\subsection{Task-Specific Resonance Windows} % [LOCKED]

Individual tasks exhibit condition-dependent resonance windows:
\begin{itemize}
    \item \texttt{nest\_005} (answer = 8360): solvable at 1-tok and
    2-tok but frozen at 3-tok and 8-tok
    \item \texttt{nest\_008} (answer = 7278): solvable \emph{only}
    at 2-tok (the sole task uniquely rescued by a single condition)
    \item \texttt{nest\_021} (answer = 5639): solvable at 2-tok and
    8-tok but frozen at 1-tok and 3-tok
\end{itemize}
More noise can \emph{prevent} solving certain tasks while enabling
others. At $n = 10$ directions, the 2-tok condition has zero frozen
tasks, while 3-tok freezes 5 and 8-tok freezes 2.

\subsection{Shi et al.\ Comparison} % [LOCKED]
\label{sec:shi}

To compare continuous and discrete perturbation regimes, we replicate
the Shi et al.~\cite{shi2025} protocol on our task set: repeating
discrete tokens ``/'' and ``?'' as 2-token prefixes, using native
embeddings (RMS $\approx 0.020$, close to our target 0.022).

\begin{table}[h]
\centering
\begin{tabular}{lccc}
\toprule
\textbf{Condition} & \textbf{Accuracy} & \textbf{$\Delta$ vs.\ baseline} & \textbf{$n$} \\
\midrule
Baseline (no prefix) & 32.0\% & --- & 1 \\
Discrete ``/'' (2 tok) & 36.0\% & +4pp & 1 \\
Discrete ``?'' (2 tok) & 48.0\% & +16pp & 1 \\
Discrete mean & 42.0\% & +10pp & 2 \\
\textbf{Random noise (2 tok)} & \textbf{51.6\%} & \textbf{+19.6pp} & \textbf{10} \\
\bottomrule
\end{tabular}
\caption{Discrete vs.\ continuous perturbation at 2 tokens.
Continuous random noise outperforms the best discrete token by 3.6pp
and the discrete mean by 9.6pp.}
\label{tab:shi}
\end{table}

Both discrete tokens improve over baseline, consistent with Shi et
al.'s observation that punctuation prefixes can aid reasoning. However,
continuous random noise in embedding space produces a substantially
larger effect (+19.6pp vs.\ +10pp mean discrete). The gap suggests
that the continuous embedding manifold offers a richer space of
perturbations than the discrete token vocabulary. Notably, ``?''
(+16pp) outperforms ``/'' (+4pp), indicating that even within the
discrete regime, token identity matters---possibly reflecting
different semantic or positional biases in the model's embedding
space.

\subsection{Cross-Model Validation} % [LOCKED]
\label{sec:crossmodel}

To test whether the perturbation effect generalizes beyond Qwen3-4B,
we run 2-token $n = 3$ scouts on three additional models using the
same 25-task set and auto-calibrated RMS scale.

\begin{table}[h]
\centering
\small
\begin{tabular}{lcccccccc}
\toprule
\textbf{Model} & \textbf{Quant} & \textbf{Base} & \textbf{+Noise} & \textbf{$\Delta$} & \textbf{Oracle} & \textbf{Rescued} & \textbf{Rescue\%} & \textbf{Regress} \\
\midrule
Qwen3-4B ($n{=}10$) & 4-bit & 32\% & 51.6\% & +19.6pp & 100\% & 17/17 & 39\% & 0 \\
Qwen3-8B ($n{=}10$) & 8-bit & 16\% & 28.8\% & +12.8pp & 80\% & 16/21 & 20\% & 0 \\
DeepSeek-1.5B ($n{=}10$) & 4-bit & 76\% & 74.4\% & $-$1.6pp & 100\% & 6/6 & --- & 0 \\
phi-2 ($n{=}3$) & none & 12\% & 18.7\% & +6.7pp & 28\% & 4/22 & 6\% & 0 \\
\midrule
Qwen3-1.7B ($n{=}3$) & 4-bit & 28\% & 29.3\% & +1.3pp & 44\% & 6/18 & --- & 2 \\
Qwen3-8B ($n{=}3$) & 4-bit & 24\% & 25.3\% & +1.3pp & 44\% & 7/19 & --- & 2 \\
\bottomrule
\end{tabular}
\caption{Cross-model 2-token perturbation results. Top: positive models
(above midrule). Bottom: null results. \emph{Rescued}: baseline
failures solved by at least one direction (out of total baseline failures).
\emph{Rescue\%}: per-direction rescue rate on baseline failures.
\emph{Regress}: baseline successes lost under all directions.
The Qwen3-8B rows constitute a within-model quantization control
(same architecture, tokenizer, task set; only quantization changed).}
\label{tab:crossmodel}
\end{table}

The effect replicates as a positive mean gain on three models: Qwen3-4B
(+19.6pp), Qwen3-8B at 8-bit quantization (+12.8pp at $n{=}10$,
oracle McNemar $p < 0.001$, 16 gains/0 losses),
and phi-2 (+6.7pp, $p = 0.125$), an out-of-family model with no thinking
mode. DeepSeek-R1-Distill-1.5B shows a distinct pattern: the mean across
10 random directions falls slightly below baseline ($-1.6$pp) but the
oracle remains 100\% (McNemar 6 gains, 0 losses, $p = 0.031$), indicating
that every baseline failure is recoverable under at least one perturbation
direction even when the average direction is unhelpful. This dissociation
between oracle accessibility and mean effect is consistent with
DeepSeek's high baseline (76\%) leaving little room for uniform gain.
The effect does not replicate on Qwen3-1.7B or Qwen3-8B at 4-bit
quantization (both +1.3pp, oracle 44\%, 2 regressions each).

\paragraph{Quantization $\times$ perturbation interaction.}
The Qwen3-8B 4-bit null is \emph{explained} by an 8-bit adjudication
experiment. Switching from 4-bit to 8-bit quantization on the same model
with identical experimental parameters produces a crossover interaction:
the baseline \emph{drops} ($24\% \to 16\%$) but the perturbation benefit
\emph{increases sharply} ($+1.3\text{pp} \to +12.8\text{pp}$ at $n{=}10$),
oracle coverage rises ($44\% \to 80\%$), and regressions vanish ($2 \to 0$).
Only 2 of 25 baseline tasks are shared between quantization levels;
oracle sets overlap on only 9 of 25.
This constitutes the cleanest within-model control in our study:
architecture, tokenizer, task set, and perturbation budget are fixed;
only weight precision varies.
The mechanism is consistent with 4-bit quantization regularizing
the trajectory landscape, helping the default reasoning path under
a fixed decode budget but washing out small perturbation signals.
8-bit preserves a richer local trajectory structure, making
perturbation-sensitive reasoning pathways accessible.

\paragraph{Substrate dependence.}
DeepSeek is \emph{smaller} than both null Qwen models (1.5B
vs.\ 1.7B and 8B at 4-bit) yet achieves 100\% oracle coverage
from $n{=}10$ directions. This rules out raw parameter count as the
predictor. The oracle/mean dissociation on DeepSeek---every baseline
failure is recoverable, but the average direction slightly
hurts---suggests that reasoning distillation creates perturbation-sensitive
trajectories that are direction-selective rather than uniformly beneficial.
This points to training regime as a key substrate requirement, alongside
sufficient quantization fidelity.

\paragraph{Cross-model dose-response.}
A dose-response on DeepSeek-R1-Distill-1.5B ($n{=}3$ per condition)
confirms a non-monotonic constructive window: 1-token perturbation
\emph{hurts} ($76\% \to 64\%$, $-12\text{pp}$), 2-token yields
the highest mean at $n{=}3$ ($+5.3\text{pp}$, SD = 0.046), and
3-token maintains a similar mean ($+4.0\text{pp}$) but enters a
bifurcated regime (SD = 0.174, Cochran's $Q = 9.5$, $p \approx
0.009$). At $n{=}10$, the 2-token mean drops to $-1.6$pp while
the oracle remains at 100\% and Cochran's $Q$ becomes significant
($p = 0.025$), indicating that additional random directions enter
the same heterogeneous regime observed at 3 tokens. This mirrors the
Qwen3-4B pattern (1-tok marginal, 2-tok peak, 3-tok decline) in
direction, establishing the non-monotonic window shape as a
cross-model phenomenon, while the oracle/mean dissociation at $n{=}10$
demonstrates that high-baseline models may benefit primarily through
task-selective direction search rather than uniform perturbation.

\paragraph{Task-level structure.} One task (\texttt{nest\_009}:
$(47 \times 92 + 50) \bmod 28 = 6$) is unlocked by all tested models
including Qwen3-1.7B, suggesting that modular arithmetic creates a
``cheap compressed trajectory'' that perturbation can discover
regardless of model capacity. Cross-model difficulty analysis reveals a
monotonic relationship: tasks that fewer models solve at baseline benefit
\emph{more} from perturbation (100\% unlock rate for tasks failed by
all models, vs.\ 0\% for universally-solved tasks).

\paragraph{Statistical significance.} Exact McNemar tests on paired
baseline-vs-oracle task outcomes: Qwen3-4B $p < 0.0001$ (17 gains, 0
losses), Qwen3-8B 8-bit $p < 0.001$ (16/0 at $n{=}10$),
DeepSeek $p = 0.031$ (6/0), phi-2 $p = 0.125$ (4/0).
Fisher's combined test across four positive models yields $p < 0.001$.

\subsection{Cross-Domain Validation: Complex Planning Tasks}
\label{sec:planning}

To test whether the perturbation effect extends beyond arithmetic, we
run a 3-way comparison on five complex planning tasks: fraud detection
system design, incident response planning, healthcare data platform
architecture, Redis cache debugging, and Oracle-to-PostgreSQL database
migration strategy. All conditions use Qwen3-4B (4-bit), greedy decoding
(temperature = 0), and \texttt{max\_new\_tokens = 2048}.

\paragraph{Conditions.}
(1)~Greedy baseline: standard generation, no modification (deterministic,
single output).
(2)~Random perturbation: 2-token random embedding-scale noise, 5 seeds.
(3)~Evolution: trained latent scorer + evolutionary search $\to$ best
evolved latent decoded via soft prompt projection, 5 seeds.

\paragraph{Finding 1: Attention sink avoidance.}
On the cache debugging task, greedy baseline produces only \textbf{14 words}
before collapsing into degenerate output---a catastrophic generation failure.
All five random perturbation seeds produce complete diagnostic plans of
\textbf{650--710 words} with systematic investigation steps, root cause
hypotheses, and resolution strategies. This binary rescue (complete failure
$\to$ complete plan) demonstrates that soft prompt perturbation in the
attention sink positions breaks degenerate greedy generation paths induced
by attention concentration patterns~\cite{xiao2024}.

\paragraph{Finding 2: Latent knowledge access via evolution.}
Evolved soft prompts do not merely produce more words---they access
qualitatively different knowledge and reasoning paths. On the incident
response task, evolution surfaces: honeypot deployment for attacker
monitoring, MITRE ATT\&CK framework analysis, tiered credential rotation
with HSM integration, immutable container rebuilds from verified base
images, and DMZ isolation of compromised services. None of these concepts
appear in baseline or random perturbation outputs. On the database
migration task, evolution converges to a structured recommendation with
regulatory-audit-aware timeline, while baseline and perturbation outputs
truncate before reaching a conclusion.

\paragraph{Caveats.}
The latent scorer is minimally trained. Evolution results are inconsistent
across seeds---some seeds surface strong latents, others do not. An
independent review ranked conditions as BASELINE $>$ EVOLUTION $>$
PERTURBATION on output ``cleanliness'' (baseline is shorter, less noisy).
The quality advantage of perturbation and evolution is in completeness
and knowledge diversity, not polish. These are five tasks with five seeds;
the findings are suggestive, not statistically powered. Better scorers,
evolution strategies, and aggregation methods should make the effect more
consistent.

\paragraph{Implications.}
The planning results suggest that embedding-space perturbation operates
through (at least) two mechanisms: (1)~breaking degenerate greedy
generation induced by attention sinks (achievable with random perturbation),
and (2)~steering attention toward different regions of the model's knowledge
space (achievable with directed evolution). The second mechanism represents
a \emph{new axis of improvement} orthogonal to scaling, fine-tuning, prompt
engineering, retrieval augmentation, and sampling-based ensembles: it changes
the continuous input representation to unlock latent knowledge with frozen
weights. The efficiency advantage over best-of-$N$ sampling is that
evolution requires $N$ forward passes through a small scorer MLP plus one
generation pass, versus $N$ full autoregressive generation passes.

\subsection{Cross-Domain Validation: Legal Reasoning Tasks}
\label{sec:legal}

To stress-test the perturbation mechanism on expert-domain reasoning, we
run a 3-way comparison on twelve complex legal reasoning scenarios spanning
five categories: regulatory framework analysis (FTC unfairness, GDPR, employment
discrimination), transactional review (SaaS contracts, startup acquisition),
risk assessment (IP portfolio, data breach triage), strategic analysis
(negotiation leverage, regulatory response), and scenario analysis (contractor
misclassification, corporate liability, whistleblower retaliation). All conditions
use Qwen3-4B (4-bit NF4), greedy decoding (temperature~=~0), and
\texttt{max\_new\_tokens~=~2048}. Evaluation is by blind Codex CLI review on
five dimensions (legal accuracy, analytical depth, practical utility, structural
quality, completeness) scored 1--10.

\paragraph{Conditions.} Identical to Section~\ref{sec:planning}: greedy
baseline (A), random perturbation with 5 seeds (B1--B5), and evolution
with trained scorer + evolutionary search and 5 seeds (C1--C5).

\paragraph{Finding 1: Oracle perturbation consistently beats baseline.}
Across all 12 tasks, the best-of-5 perturbation output
scores higher than greedy baseline in \textbf{11 out of 12 tasks} (92\%).
Oracle lifts range from $+0.6$ to $+3.4$ on the 10-point expert-review
scale, with an average lift of $+1.6$ points. The sole exception is a
GDPR controller/processor task where the model lacks sufficient domain
knowledge regardless of condition.

\paragraph{Finding 2: Perturbation unlocks qualitatively different legal analysis.}
On task 10 (contractor misclassification), the greedy baseline misstates
the California ABC test and produces catastrophic math errors (score:
2.2/10). The best perturbation seed provides a full cross-jurisdiction
analysis across California, New York, Texas, and federal law with
triage priorities and risk quantification (score: 5.6/10). On task 08
(negotiation leverage), baseline produces a badly truncated 379-word
fragment (2.0/10), while perturbation produces complete negotiation
strategies with demand-by-demand analysis (5.4/10). These are not
incremental improvements---perturbation accesses legal reasoning
trajectories the model cannot reach via greedy decoding.

\paragraph{Finding 3: Evolution proves the concept, judges determine the outcome.}
On the one task with fully functional evolution (task 01, FTC unfairness),
evolution produces the highest group average (6.2/10 vs.\ perturbation
5.8 vs.\ baseline 5.2). Evolution outputs contain zero fabricated
citations and zero internal contradictions, versus 2/5 contradictions in
perturbation outputs. However, evolution was impaired on 9/12 tasks due
to a scorer initialization bug (non-deterministic random projection layer),
producing degenerate identical outputs. When evolution works, the trained
scorer guides latent search toward higher-quality trajectories. When it
fails, the results clearly demonstrate that \emph{scorer quality is the
binding constraint}---the latent knowledge exists (proven by oracle
perturbation), and the judge's ability to identify good trajectories
determines whether that knowledge is reliably accessed.

\paragraph{Implications for judge-heavy systems.}
These results establish that embedding-space perturbation is a
\emph{mechanism} for accessing latent knowledge, while the scorer/judge
is the \emph{selector} that determines which trajectory is surfaced to
the user. The oracle analysis proves the ceiling is high (92\% of tasks
improvable). The gap between oracle performance and mean performance
represents the opportunity for better judges. This is by design: the
system is explicitly judge-heavy, and judge quality determines outcome
quality. A production-quality trained judge (as opposed to the minimally
trained scorer used here) should capture substantially more of the
demonstrated oracle ceiling.

% ============================================================================
\section{Discussion}

\subsection{Two-Level Trajectory Model}

The structural decomposition suggests a two-stage model:
\begin{enumerate}
    \item \textbf{Magnitude sets the task-level threshold.}
    Perturbation energy determines which tasks are \emph{unanimous}
    (solved by all directions), \emph{frozen} (solved by none), or
    \emph{sensitive} (solved by some). At $n = 10$, two tokens freeze
    zero tasks, three tokens freeze 5, and eight tokens freeze 2.
    The three-token frozen set is a strict superset of the two-token
    frozen set: over-perturbation contracts reachability.

    \item \textbf{Direction selects within the sensitive set.}
    Each random direction solves a different subset of the sensitive
    tasks. Per-direction solve counts vary normally across all
    conditions (all $p > 0.3$ vs.\ heterogeneous iid null), so the
    diversity is in task identity, not in total solve count.
\end{enumerate}

The two-token optimum emerges because it maximizes the reachable task
set (25/25 at $n = 10$, vs.\ 20/25 for three tokens) while maintaining
the highest mean per-direction accuracy (12.9/25 = 51.6\%).

\subsection{Convergence vs.\ Computation}

A grading-confound analysis reveals that the accuracy effect has
model-dependent character. Extracting all integers from each response
and checking whether the correct answer appears anywhere (not just as
the last integer) shows:

\begin{table}[h]
\centering
\begin{tabular}{llccc}
\toprule
\textbf{Model} & \textbf{Grading Rule} & \textbf{Baseline} & \textbf{Perturbed} & \textbf{$\Delta$} \\
\midrule
Qwen3-4B & Last integer (used) & 32\% & 51.6\% & +19.6pp \\
Qwen3-4B & Answer anywhere & 80\% & 82\% & +2pp \\
\midrule
Qwen3-8B (8-bit) & Last integer (used) & 16\% & 22\% & +6pp \\
Qwen3-8B (8-bit) & Answer anywhere & 32\% & 50\% & +18pp \\
\bottomrule
\end{tabular}
\caption{Accuracy under two grading rules. Qwen3-4B can already compute
the correct answer 80\% of the time; its gain is purely convergence
(placing the answer last). Qwen3-8B at 8-bit has a lower computational
ceiling (32\%); perturbation appears to improve computation as well as
convergence. \emph{Caveat}: answer-anywhere is computed from stored
traces truncated at 2000 characters ($>$80\% of responses are
truncated), so these rates are lower bounds.}
\label{tab:convergence}
\end{table}

For Qwen3-4B, the model can already derive the correct answer in its
reasoning chain for 80\% of tasks but fails to produce it as the
final output. Perturbation improves convergence without changing
computational success. For Qwen3-8B at 8-bit, the model has substantial
room for computational improvement (32\% answer-anywhere at baseline),
and perturbation produces apparent gains in both computation (+18pp)
and convergence (+6pp), though these estimates are based on truncated
stored traces and should be treated as preliminary. This pattern is
consistent with the mechanism depending on the model's computational
ceiling, but confirming this requires full-trace analysis.

\subsection{Trajectory Efficiency Examples}

Response inspection confirms the convergence mechanism concretely:
\begin{itemize}
    \item \texttt{nest\_010} ($11 \times 21 \times 4 = 924$): The
    baseline response enters a verbose place-value explanation that
    exhausts the token budget before reaching the answer. Under noise,
    the model takes a more direct path ($11 \times 21 = 231$,
    $231 \times 4 = 924$) and completes within budget.
    \item \texttt{nest\_006} ($= 6$): The baseline derives the correct
    answer within its reasoning trace but continues generating,
    runs out of tokens, and the last integer is not 6. Under perturbation,
    the model produces a clean answer ending with $\boxed{6}$---a
    convergence fix, not a computation fix.
\end{itemize}
The effect is trajectory efficiency: noise routes computation through
paths that terminate at the correct answer within the generation budget.

\subsection{Why Two Tokens?}

At $n = 10$ directions, two tokens have zero frozen tasks, compared to
5 at three tokens and 2 at eight tokens. At matched budget ($k = 3$),
two tokens freeze 5.8 tasks on average versus 9.9 for three tokens
and 7.9 for eight tokens. Two tokens thus maximize the reachable task
set while maintaining the highest per-direction accuracy
(mean 12.9/25 = 51.6\%). We hypothesize that two tokens provide the
minimum perturbation energy needed to productively diversify the
reasoning trajectory while staying below the energy level that
freezes marginal tasks. At one token, the perturbation is marginal;
at three or more, the model's reasoning trajectory diverges too early,
exhausting the token budget before reaching an answer.

\subsection{Implications for Practice}

The training-free nature of this intervention suggests a practical
``perturbation ensemble'' strategy: run a small model with $k$ random
two-token prefixes and take the majority vote or oracle. Ten two-token
directions achieve 100\% coverage (25/25 tasks), while ten three-token
directions reach only 80\% and add no tasks beyond the two-token
oracle. Beyond the two-token optimum, additional perturbation is
deleterious: it lowers per-direction accuracy and contracts the
reachable task set, freezing tasks that are solvable at two tokens.

% ============================================================================
\section{Limitations} % [LOCKED]
\label{sec:limitations}

\begin{itemize}
    \item \textbf{Model-contingent.} Cross-model scouts
    (Section~\ref{sec:crossmodel}) show positive transfer to four models
    (Qwen3-4B, Qwen3-8B at 8-bit, DeepSeek-1.5B, phi-2) but not to
    Qwen3-1.7B, suggesting the effect is substrate-dependent rather than
    universal. Full characterization across architectures and scales
    remains future work.

    \item \textbf{Cross-domain evidence growing but scorer-limited.}
    Planning task validation (Section~\ref{sec:planning}) demonstrates
    qualitative effects on five tasks. Legal reasoning validation
    (Section~\ref{sec:legal}) provides stronger cross-domain evidence:
    oracle perturbation beats baseline in 11/12 tasks with average
    $+1.6$ lift. However, evolution results depend on a barely-trained
    scorer that was impaired on 9/12 tasks due to a projection layer bug.
    The system is explicitly judge-heavy: the perturbation mechanism
    reliably accesses latent knowledge (proven by oracle analysis), but
    the degree to which that knowledge is captured depends entirely on
    judge/scorer quality. Better-trained judges should narrow the gap
    between oracle and mean performance.

    \item \textbf{Small $n$.} 25 tasks with 10 perturbation directions at
    the key 2-token condition. Per-task mechanistic analyses remain
    statistically limited.

    \item \textbf{Quantization--perturbation entanglement.}
    The Qwen3-8B within-model control (Section~\ref{sec:crossmodel})
    demonstrates that quantization modulates perturbation sensitivity:
    4-bit yields a null result while 8-bit yields +16pp.
    This raises the question of whether our primary Qwen3-4B results
    are quantization-specific. Full-precision experiments are needed
    to determine whether 4-bit represents a sweet spot or whether
    higher fidelity would yield even larger effects.

    \item \textbf{Token budget interaction.} Regressions on easy tasks
    (large-answer category: 100\% $\to$ 78\%) and the 8-bit 8B baseline
    degradation ($16\%$ vs.\ $24\%$ at 4-bit, with 8-bit traces being
    more verbose) suggest the effect interacts with the 1024-token
    generation budget. Perturbation may benefit models that can reach
    correct answers within budget, while harming those whose reasoning
    chains are already budget-limited.

    \item \textbf{Convergence vs.\ computation confound.}
    A grading-confound analysis (Table~\ref{tab:convergence}) reveals
    the model can already compute correct answers 80\% of the time
    at baseline. For Qwen3-4B, the accuracy gain is primarily a
    convergence effect (placing the correct answer last). For
    Qwen3-8B at 8-bit, perturbation also improves computation
    (answer-anywhere 32\% $\to$ 50\%). The mechanism appears to
    depend on the model's computational ceiling.
    DeepSeek shows the reverse: perturbation degrades both
    answer-anywhere (84\% $\to$ 78\%) and last-integer accuracy
    (76\% $\to$ 69\%), with only oracle coverage benefiting.
\end{itemize}

% ============================================================================
\section{Conclusion}

Random prefix perturbations at the embedding scale diversify reasoning
trajectories in small language models, producing a sharp non-monotonic
optimum. Two tokens raise Qwen3-4B mean accuracy from 32\% to 51.6\%.
A grading-confound analysis reveals the mechanism is model-dependent:
for Qwen3-4B (where the model already computes correct answers 80\%
of the time), the gain is primarily \emph{convergence}---placing the
correct answer as the final output. For Qwen3-8B at 8-bit (lower
computational ceiling at 32\% answer-anywhere), perturbation genuinely
improves computation (+18pp answer-anywhere) as well as convergence.

The trajectory diversification is genuine and well-characterized:
different random directions solve different task subsets, enabling
100\% oracle coverage from 10 two-token directions on Qwen3-4B and
80\% on Qwen3-8B at 8-bit. Per-direction solve counts vary normally;
the diversity is in task identity, not total count.
Cross-model validation confirms the effect on four models, while
a within-model quantization control reveals that weight precision
modulates perturbation sensitivity (4-bit null, 8-bit +12.8pp at
$n{=}10$). A cross-model dose-response on DeepSeek confirms the
non-monotonic window generalizes: 1-token hurts, 2-token is the
local optimum, 3-token enters a bifurcated regime. At $n{=}10$,
DeepSeek exhibits oracle/mean dissociation---every baseline failure
is recoverable but the average direction slightly hurts.

Cross-domain validation on five complex planning tasks reveals two
additional phenomena. First, perturbation rescues catastrophic
attention-sink-induced generation failures: a 14-word baseline
collapse becomes a 650+ word complete plan across all five perturbation
seeds. Second, evolved soft prompts (via a trained latent scorer +
evolutionary search) surface qualitatively different reasoning---security
frameworks, architectural patterns, and diagnostic strategies absent
from both baseline and random perturbation outputs---suggesting that
directed embedding-space perturbation can access latent knowledge the
model possesses but cannot reach under standard decoding.

Third, on twelve complex legal reasoning tasks (Section~\ref{sec:legal}),
the best-of-5 perturbation output beats greedy baseline in 11/12
tasks with an average lift of $+1.6$ on a 10-point blind
expert-review scale. The legal domain confirms the mechanism's
generality: perturbation unlocks multi-jurisdictional reasoning,
regulatory frameworks, and strategic analysis that greedy decoding
cannot access. Evolution with a minimally-trained scorer wins outright
on the one fully-functional task, establishing that the system is
explicitly \emph{judge-heavy}---latent knowledge access is reliable,
and judge quality determines how much of the oracle ceiling is captured.

These findings establish embedding-space perturbation as a training-free
intervention that diversifies reasoning trajectories and, when directed
by evolution, can unlock latent knowledge inaccessible under standard
decoding. Accuracy gains in arithmetic reflect a mixture of improved
convergence and computation depending on the model's baseline capability;
in planning, the mechanism breaks degenerate generation paths and steers
attention toward different knowledge regions; in legal reasoning, it
accesses domain-specific analysis the model possesses but cannot reach
via greedy decoding. The gap between oracle performance (92\% of tasks
improvable) and mean performance quantifies the opportunity for better
judges---the binding constraint in this judge-heavy architecture.

% ============================================================================
\bibliographystyle{plain}
% References (inline for now; move to .bib for submission)
\begin{thebibliography}{99}

\bibitem{shi2025}
Shi, Z., et al.
\newblock Meaningless tokens, meaningful gains: Repeated tokens improve LLM reasoning.
\newblock \emph{arXiv:2510.01032}, 2025.

\bibitem{goyal2024}
Goyal, S., et al.
\newblock Think before you speak: Training language models with pause tokens.
\newblock \emph{ICLR}, 2024.

\bibitem{salinas2024}
Salinas, D. and Morstatter, F.
\newblock The butterfly effect of altering prompts: How small changes and jailbreaks affect large language model performance.
\newblock \emph{arXiv:2401.03729}, 2024.

\bibitem{london2025}
London, B. and Nagarajan, V.
\newblock Provably more expressive transformers via extra tokens.
\newblock \emph{NeurIPS}, 2025.

\bibitem{li2025}
Li, Z., et al.
\newblock Quasi-Lyapunov exponent for language models.
\newblock \emph{arXiv}, 2025.

\bibitem{gammaitoni1998}
Gammaitoni, L., et al.
\newblock Stochastic resonance.
\newblock \emph{Reviews of Modern Physics}, 70(1):223, 1998.

\bibitem{qwen2025}
Qwen Team.
\newblock Qwen3 technical report.
\newblock 2025.

\bibitem{xiao2024}
Xiao, G., et al.
\newblock Efficient streaming language models with attention sinks.
\newblock \emph{ICLR}, 2024.

\bibitem{xu2025}
Xu, Z., et al.
\newblock Causal amplification in the residual stream.
\newblock \emph{arXiv}, 2025.

\bibitem{kim2026rsp}
Kim, S., et al.
\newblock From noise to diversity: Random embedding injection in LLM reasoning.
\newblock \emph{arXiv:2605.11936}, 2026.

\bibitem{zhu2025softreasoning}
Zhu, Y., et al.
\newblock Soft reasoning: scalable, model-agnostic reasoning via embedding perturbation.
\newblock \emph{arXiv:2505.24688}, 2025.

\bibitem{guo2025swift}
Guo, Z., et al.
\newblock Mining intrinsic rewards from LLM hidden states for efficient best-of-N sampling.
\newblock \emph{arXiv:2505.12225}, 2025.

\end{thebibliography}

\end{document}
